%%%%%%%%%%%%%%%%%%%%%%%%%%%%%%%%%%%%%%%%%%%%%%%%%%%%%%%%%%%%%%%%%%%%%%%

\ifdefined\whole
\else
\documentclass{amsart}
\usepackage{amstext,amsfonts,amssymb,amscd,amsbsy,amsmath,verbatim, mathrsfs, fullpage}
\usepackage[alphabetic,backrefs,abbrev,lite]{amsrefs} % for bibliography
\usepackage[colorlinks=true,hyperindex,linkcolor=magenta,pagebackref=true,citecolor=cyan]{hyperref}
\usepackage[dvipsnames]{xcolor} % for colors
\usepackage{ifthen,tikz}
\usetikzlibrary{shapes.geometric}
\usepackage{tikz-cd}
\usepackage{tikz-3dplot}
\usepackage{color}
\usepackage{amsthm}
\usepackage{cleveref}
\usepackage{latexsym}
%\usepackage[color, matrix, arrow]{xy}
\usepackage{enumerate}
\usepackage{mathtools,accents}
\usepackage[all]{xy}
%\setcounter{MaxMatrixCols}{15} -- this line was causing an ``options class with xy''
%\usepackage{pgfplots}
%\pgfplotsset{compat=1.18}


\def\AA{{\mathbb A}}
\def\sO{{\mathcal O}}
\def\de#1{{\bf *** #1 ***}}
\DeclarePairedDelimiter\ceil{\lceil}{\rceil}
\DeclarePairedDelimiter\floor{\lfloor}{\rfloor}

\numberwithin{equation}{section}

\newtheorem{lemma}[equation]{Lemma}
\newtheorem{theorem}[equation]{Theorem}
\newtheorem*{gr}{Green's Linear Syzygy Theorem}
\newtheorem{propo}[equation]{Proposition}
\newtheorem{prop}[equation]{Proposition}
\newtheorem{cor}[equation]{Corollary}
\newtheorem{conj}[equation]{Conjecture}
\newtheorem{claim}[equation]{Claim}
%\newtheorem{claim*}{Claim}
\newtheorem{thm}[equation]{Theorem}
\newtheorem{convention}[equation]{Convention}
\newtheorem{question}[equation]{Question}


\theoremstyle{definition}
\newtheorem{defn}[equation]{Definition}
\newtheorem{example}[equation]{Example}
\newtheorem{notation}[equation]{Notation}
\newtheorem{construction}[equation]{Construction}
\newtheorem{algorithm}[equation]{Algorithm}
\newtheorem{warning}[equation]{Warning}
\newtheorem{conventions}[equation]{Conventions}

\theoremstyle{remark}
\newtheorem{remark}[equation]{Remark}
\newtheorem{remarks}[equation]{Remarks}
\newtheorem{rem}[equation]{Remark}
\newtheorem*{claim*}{Claim}
\newtheorem*{case*}{Case}

\numberwithin{figure}{section}
% Commands

\newcommand{\Kos}{\mathcal{K}}
\newcommand{\cS}{\mathcal{S}}
\newcommand{\rC}{\mathrm{C}}
\newcommand{\Cech}{\check{\mathrm{C}}}
\newcommand{\Tate}{\on{Tate}}
\newcommand{\Tail}{{\mathbf{Tail}}}
\newcommand{\cC}{\mathcal{C}}
\newcommand{\cK}{\mathcal{K}}
\newcommand{\tot}{\operatorname{tot}}

\newcommand{\isom}{\cong}
\newcommand{\m}{\mathfrak m}
\newcommand{\PP}{\mathbb P}
\newcommand{\PPv}{\mathbb P_{\on{var}}}
\newcommand{\bD}{\mathbf D}
\newcommand{\df}{\operatorname{diff}}
\renewcommand{\P}{\PP}
\newcommand{\bA}{\mathbb A}
\newcommand{\A}{\bA}
\newcommand{\HH}{H}
\newcommand{\GG}{\mathbb G}
\newcommand{\ZZ}{\mathbb Z}
\newcommand{\NN}{\mathbb N}
\newcommand{\QQ}{\mathbb Q}
\newcommand{\bH}{\mathbf H}
\newcommand{\bF}{\mathbf F}
\newcommand{\DD}{\mathbf D}
\newcommand{\lideal}{\langle}
\newcommand{\rideal}{\rangle}
\newcommand{\initial}{\operatorname{in}}
\newcommand{\pdim}{\operatorname{pdim}}
\newcommand{\Hilb}{\operatorname{Hilb}}
\newcommand{\Spec}{\operatorname{Spec}}
\newcommand{\NE}{\overline{\operatorname{NE}}}
\newcommand{\Eff}{\operatorname{Eff}}
\newcommand{\im}{\operatorname{im}}
\newcommand{\NS}{\operatorname{NS}}
\newcommand{\Frac}{\operatorname{Frac}}
\newcommand{\ch}{\operatorname{char}}
\newcommand{\Proj}{\operatorname{Proj}}
\newcommand{\id}{\operatorname{id}}
\newcommand{\Div}{\operatorname{Div}}
\newcommand{\tr}{\operatorname{tr}}
\newcommand{\Tr}{\operatorname{Tr}}
\newcommand{\Supp}{\operatorname{Supp}}
\newcommand{\Gal}{\operatorname{Gal}}
\newcommand{\Pic}{\operatorname{Pic}}
\newcommand{\QQbar}{{\overline{\mathbb Q}}}
\newcommand{\Br}{\operatorname{Br}}
\newcommand{\Bl}{\operatorname{Bl}}
\newcommand{\Cox}{\operatorname{Cox}}
\newcommand{\Tor}{\operatorname{Tor}}
\newcommand{\diam}{\operatorname{diam}}
\newcommand{\Hom}{\operatorname{Hom}} %done
\newcommand{\sheafHom}{\mathcal{H}om}
\newcommand{\Gr}{\operatorname{Gr}}
\newcommand{\HF}{\operatorname{HF}}
\newcommand{\HP}{\operatorname{HP}}
\newcommand{\Osh}{{\mathcal O}}
\newcommand{\cO}{{\mathcal O}}
\newcommand{\kk}{{\bf k}}
\newcommand{\rank}{\operatorname{rank}}
\newcommand{\length}{\operatorname{length}}
\newcommand{\codim}{\operatorname{codim}}
\newcommand{\depth}{\operatorname{depth}}
%\newcommand{\FF}{\mathbb{F}}
\newcommand{\F}{\FF}
\newcommand{\Sym}{\operatorname{Sym}} %done
\newcommand{\GL}{{GL}}
\newcommand{\R}{\mathbb{R}}
\newcommand{\CC}{\mathbb{C}}
\newcommand{\Syz}{\operatorname{Syz}}
\newcommand{\Prob}{\operatorname{Prob}}
\newcommand{\defi}[1]{\textsf{#1}} % for defined terms
\newcommand{\Htot}{H_{\tot}}
\newcommand{\Ltot}{L_{\tot}}
\newcommand{\beq}{\begin{displaymath}}
\newcommand{\eeq}{\end{displaymath}}
\newcommand{\bs}{\backslash}
\newcommand{\ff}{\mathbf{f}}
\newcommand{\Gam}{\Gamma}
\newcommand{\Lotimes}{\overset{L}{\otimes}}
\newcommand{\uHom}{\underline{\Hom}}
\def\nc{\newcommand}
\nc{\tors}{\on{tors}}



\newcommand{\Bmod}{\ensuremath{B_\text{mod}}}
\newcommand{\Bint}{\ensuremath{B_\text{int}}}
\newcommand\commentr[1]{{\color{red} \sf [#1]}}
\newcommand\commentb[1]{{\color{blue} \sf [#1]}}
\newcommand\commentm[1]{{\color{magenta} \sf [#1]}}
\newcommand{\daniel}[1]{{\color{blue} \sf $\clubsuit\clubsuit\clubsuit$ Daniel: [#1]}}
\newcommand{\michael}[1]{{\color{red} \sf $\clubsuit\clubsuit\clubsuit$ Michael: [#1]}}
\newcommand{\maya}[1]{{\color{green} \sf $\clubsuit\clubsuit\clubsuit$ Maya: [#1]}}
\newcommand{\mahrud}[1]{{\color{ForestGreen} \sf $\blacklozenge$ Mahrud: [#1]}}
\newcommand{\lauren}[1]{{\color{Fuchsia} \sf $\diamondsuit$ Lauren: [#1]}}
\newcommand{\david}[1]{{\color{magenta} \sf ** David: [#1] **}}
\newcommand{\jay}[1]{{\color{teal} \sf $\spadesuit\spadesuit\spadesuit$ Jay: [#1]}}

\def\edim{\operatorname{edim}}
\def\reg{\operatorname{reg}}

\newcommand{\ve}[1]{\ensuremath{\mathbf{#1}}}
\newcommand{\chr}{\ensuremath{\operatorname{char}}}

\newcommand{\idealof}{\operatorname{ideal}}
\newcommand{\Nef}{\operatorname{Nef}}

%Added by MB:
\def\nc{\newcommand}
\def\on{\operatorname}
\nc{\Q}{\mathbb{Q}}
\nc{\RR}{\mathbf{R}}
\nc{\LL}{\mathbf{L}}
\nc{\cH}{\mathcal H}
\nc{\xra}{\xrightarrow}
\nc{\xla}{\xleftarrow}
\def\a{\mathbf{a}}
\def\om{\omega}
\def\Om{\Omega}
\def\DM{\operatorname{DM}}
\def\Coh{\operatorname{Coh}}
\def\Mod{\operatorname{Mod}}
\def\free{\operatorname{free}}
\def\QCoh{\operatorname{QCoh}}
\def\Cpx{\operatorname{Cpx}}
\def\th{\on{th}}
\def\F{\mathcal{F}}
\def\coker{\on{coker}}
\def\p{\partial}
\def\wt{\widetilde}
\def\i{\iota}
\nc{\into}{\hookrightarrow}
\nc{\onto}{\twoheadrightarrow}
\nc{\OO}{\mathcal{O}}
\nc{\Z}{\mathbb{Z}}
\nc{\cA}{\mathcal{A}}
\nc{\w}{\widehat}
\nc{\End}{\on{End}}
\nc{\res}{\frac{1}{x_0x_1}}
\nc{\tF}{\widetilde{F}}
\nc{\tG}{\widetilde{G}}
\nc{\tf}{\widetilde{f}}
\nc{\Com}{\on{Com}}
\nc{\exact}{\on{exact}}

\nc{\G}{\mathbb{G}}
\nc{\cG}{\mathcal{G}}
\nc{\cE}{\mathcal E}
\nc{\cF}{\mathcal F}
\nc{\cR}{\mathcal R}
\nc{\cD}{\mathcal D}
\nc{\cB}{\mathcal B}
\nc{\cT}{\mathcal T}
\nc{\cL}{\mathcal L}
\def\M{\mathcal{M}}
\nc{\bM}{\mathbf M}
\nc{\bN}{\mathbf N}
\nc{\U}{\mathbf U}
\nc{\BM}{\mathbf B \mathbf M}
\nc{\Dsg}{\on{D}_{\on{sg}}}
\nc{\fC}{\mathcal{C}}
\nc{\fG}{\mathcal{G}}
\nc{\N}{\mathbb{N}}



%When merging files, add these
\nc{\del}{\partial}
\nc{\cone}{\on{cone}}
\nc{\D}{\on{D}_{\on{diff}}}
\nc{\DMb}{\on{D}^b_{\DM}}
\nc{\Db}{D}
\nc{\Kb}{\on{K}^{\on{b}}}
\nc{\fm}{\mathfrak{m}}
\nc{\Flag}{\on{Flag}}
\nc{\DMmin}{\DM_{\on{min}}}
\nc{\Ddiff}{\on{D}_{\on{diff}}}
\nc{\Dbdiff}{\on{D}^\on{b}_{\on{diff}}}
\nc{\wO}{\widehat{\OO}}
\nc{\wT}{\widehat{T}}
\nc{\from}{\leftarrow}
\nc{\wLL}{\widetilde{\LL}}
\nc{\augCech}{\widetilde{\cC}}
\nc{\Fold}{\on{Fold}}
\nc{\Ext}{\on{Ext}}
\nc{\RHom}{\on{RHom}}
\nc{\FF}{\mathbf{F}}
\nc{\Comper}{\Com_{\on{per}}}
\nc{\Unfold}{\on{Unfold}}
\nc{\intHom}{\underline{\Hom}}
\nc{\Ex}{\on{Ex}}
\nc{\tg}{\widetilde{g}}
\def\tP{\widetilde{P}}
\def\b{\beta}
\nc{\B}{\mathcal{B}}
\nc{\K}{\mathcal{K}}
\nc{\kos}{\on{Kos}}
\nc{\Perf}{\on{Perf}}
\nc{\tR}{\widetilde{\cR}}
\nc{\X}{\mathcal{X}}
\nc{\Cl}{\on{Cl}}
\nc{\fU}{\mathcal{U}}
\nc{\bU}{\mathbf U}
\def\co{\colon}
\def\ce{\coloneqq}
\nc{\st}{\on{st}}
\def\E{\mathcal{E}}
\nc{\coh}{\on{coh}}
\def\ex{\on{ex}}
\def\D{D}
\def\lin{\on{lin}}
\def\g{\gamma}
\nc{\tU}{\U}
\nc{\bC}{\mathbf{C}}
\nc{\aux}{\on{aux}}
\def\d{\mathbf{d}}
\def\phi{\varphi}
\def\cP{\mathcal{P}}
\def\I{\mathcal{I}}
\def\geo{\on{geo}}
\def\T{\mathbf{T}}
\nc\Dsing{\on{D}^{\on{sing}}}
\def\Y{\mathcal{Y}}
\def\bL{\mathbf{L}}
\def\k{\mathbf{k}}
\def\e{\epsilon}
\def\sX{\mathcal X}
\def\tsX{\widetilde{\mathcal X}}
\def\epsilon{\varepsilon}
\def\w{\mathbf w}
\def\colim{\on{colim}}
\def\fix#1{{{\bf ***}#1{\bf ***}}}
\def\xla{\xleftarrow}
%\begin{document}
%\title[Toric methods in homological algebra]{Arbeitsgemeinschaft:  \\ Toric methods in homological algebra \\ October 11 - 16, 2026}
%\title{Section from Notes for the Arbeitsgemeinschaft}
%\author{Michael K. Brown}
%\author{David Eisenbud}
%\author{Daniel Erman}
%This command gets rid of the MR number in the bibliography
\def\MR#1{}

\newcommand{\xto} {\xrightarrow}
\newcommand{\xgets}{\xleftarrow}

\begin{document}
\fi

\section{Good multigradings}\label{sec:good-gradings}

Let $S=k[x_1, \dots, x_n]$ be a polynomial ring graded by an arbitrary finitely generated abelian group $\Cl(S)$ via a map $\ZZ^n\overset{\deg}{\longrightarrow}\Cl(S)$, where $\epsilon_1, \dots, \epsilon_n$ are a fixed basis of $\ZZ^n$ and we set $\deg(x_i) \ce\deg(\epsilon_i)\in \Cl(S)$.  This induces a grading on $S$ where $\deg(x_1^{b_1}\cdots x_n^{b_n})= \deg b$ for $b\in\ZZ^n$ and $\deg(a)=0$ for $a\in k$.  We will assume that the map $\deg$ is surjective.

Our goal in this section is to articulate conditions under which the multigraded ring $S$ corresponds to the Cox ring of some projective toric variety.  We will need some notation.  Write $M\ce\ker(\deg)$ and consider the induced short exact sequence:
\[
0\to M\overset{\alpha}{\to} \ZZ^n \overset{\deg}{\longrightarrow} \Cl(S) \to 0.
\]
For convenience, we fix bases of $M$ and $\ZZ^n$ so that we can represent $\alpha$ as a matrix.  Define $\Eff(S) \ce\{d \in \Cl(S) \mid S_d \ne 0\}$, the \defi{effective semigroup}.  We can consider the corresponding cone $\Eff(S)_{\QQ}\subseteq \Cl(S)\otimes_{\ZZ}\QQ$, the positive rational span of the image of $\Eff(S)$ under the natural map $\Cl(S)\otimes_{\ZZ}\QQ$.

\begin{defn}\label{defn:posGrading}
  We say that $\deg\colon\ZZ^n\to\Cl(S)$ is a \defi{positive grading} if there exists a group homomorphism $\phi\colon \Cl(S) \to \ZZ$ such that $\phi \circ \deg(\epsilon_i) > 0 $ for all $i$.
\end{defn}
\begin{lemma}\label{lem:positive-gradings}
  The following are equivalent:
  \begin{enumerate}
    \item $S$ is positively graded.
    \item $S_0=k$.
    \item $\Eff(S)_{\QQ}$ is a pointed cone.
  \end{enumerate}
\end{lemma}
\begin{proof}
    $(1)\Rightarrow (2)$:  By assumption $k$ lies in $S_0$.  If another polynomial lies in $S_0$ then there exists $b\neq 0$ such that $\deg b=0$.  Thus $\sum b_i\deg\epsilon_i=0$ so $\sum b_i\varphi\circ\deg(\epsilon_i)=0$ for some $\varphi$ witnessing the positivity of the grading, which is impossible.

    $(2)\Rightarrow (3)$:  If $\Eff(S)_{\QQ}$ is not a pointed cone then there exists $d\in \Cl(S)$ with $d\neq 0$ where $d$ and $-d$ both lie in $\Eff(S)$.  Choose $f\in S_d$ and $g\in S_{-d}$, so that $\deg(fg)=0$.  Since $S_0=k$ both $f$ and $g$ are units and thus scalars, contradicting $d\neq 0$.

    $(3)\Rightarrow (1)$:  Define the positive grading $\phi$ by choosing a projection of vector spaces $\Cl(S)\otimes \QQ\to \QQ$ that sends $\Eff(S)_{\QQ}\setminus \{0\}$ into $\QQ_{>0}$.  This is possible because $\Eff(S)_{\QQ}$ is a pointed cone. Then clear denominators as needed so that the image of $\Eff(S)$ lies in $\ZZ$.
  \end{proof}

\begin{prop}\label{prop:goodMultigradings}
The two conditions below are equivalent for $S$ and $\deg$ as above (and are not affected by a change of basis for $M$):
\begin{enumerate}
    \item The ring $S$ and degree map $\deg$ satisfy the following:
      \begin{enumerate}
        \item The map $\deg$ is surjective;
        \item
        Each row of $\alpha$ is the primitive point on a distinct ray, meaning that the rows of
        $\alpha$ are distinct and nonzero, and the entries of each row have greatest common divisor 1;
        \item The grading on $S$ defined by $\deg$ is a positive grading.
    \end{enumerate}
    \item  There exists a projective, simplicial toric variety $X$ \lauren{with no torus factors??} whose Cox ring is $S$ as a multigraded ring, in the following specific sense: The fan for $X$ contains $n$ rays and we have an isomorphism $\Cl(X)\cong \Cl(S)$ such that for each $i$ the image $\deg(x_i)\in \Cl(S)$ is the class in $\Cl(X)$ of the torus invariant divisor corresponding to the $i$th ray of $X$.
\end{enumerate}
\end{prop}


  % \begin{prop}\label{prop:goodMultigradings}
  %   The following are equivalent: \lauren{the logic of this reads very strangely to me currently .. there does indeed exist a polynomial ring satisfying the conditions in (1) so it seems like (2) should always be true}\daniel{Good point, totally agree.  Rewrite as you see fit.}
  %   \begin{enumerate}
  %     \item There is a polynomial ring $S$, graded via $\deg\colon \ZZ^n\to \Cl(S)$, that satisfies the following conditions:
  %     \begin{enumerate}
  %       \item $\deg$ is surjective;
  %       \item $S$ is positively graded;
  %       \item Each row of $\alpha$ \lauren{what is $\alpha$?}\daniel{It was a name I gave to the matrix $M\to \ZZ^n$.} is the primitive point on a distinct ray; that is, the rows of
  %       $\alpha$ are distinct, nonzero, and {\color{blue} the entries of each row} have greatest common divisor 1. \lauren{what does it mean for vectors to have greatest common divisor 1?}
  %       %matrix representing $\alpha\colon M\to \ZZ^n$ has no zero rows and no rows that are scalar multiples of one another.
  %     \end{enumerate}
  %     \item There is a simplicial, projective toric variety $X$ with $\Cl(X) = \Cl(S)$ and where the Cox ring of $X$ equals $S$.
  %   \end{enumerate}
  % \end{prop}
  \begin{proof}
    $(2)\Rightarrow (1)$ follows from standard facts about toric varieties.

    Given such an $X$, the degree map $\deg\colon\ZZ^n\to\Cl(X)$ is surjective because the classes of torus-invariant divisors generate $\Cl(X)$ \cite[Theorem 4.1.3]{CLS2011}.  Choose a basis for the lattice $N$ containing the fan for $X$ and take the dual basis for the lattice $M$, giving a matrix $N\gets \ZZ^n$ containing primitive elements along each ray of the fan.  Since the rays are distinct and nonzero the columns of this matrix are nonzero and no column is a scalar multiple of another. Dualizing then yields a map $\alpha\colon M\to \ZZ^n$ satisfying condition (b).  The grading $\deg$ is positive by \cite[Corollary 1.2(i)]{Cox95} and Lemma~\ref{lem:positive-gradings}.

    $(1)\Rightarrow (2)$ is more subtle.  We will be using facts from~\cite[Chapters 14 and 15]{CLS2011}.

    The beginning of \cite[Chapter 14.4]{CLS2011} defines a secondary fan for arbitrary algebraic actions of a subgroup $G\subseteq (k^*)^n$ on $\Spec(S)\cong \mathbb A^n$ (see for instance \cite[Lemma 14.4.6,7]{CLS2011}).  Our map $\deg\colon\ZZ^n\to\Cl(S)$ defines a map of group rings $k[\ZZ^n]\to k[\Cl(S)]$.  The surjectivity in (a) implies that the corresponding map of algebraic groups defines a subgroup of $\Spec k[\ZZ^n]\cong(k^*)^n$. \lauren{is this right?} The \lauren{possibly nonexistent?} chambers of the secondary fan for the induced action of this subgroup on $\Spec S$ correspond to GIT quotients for different choices of linearized line bundle.  We will show that at least one has the properties we desire.

    In particular, \cite[Proposition 15.1.6]{CLS2011} describes when the secondary fan contains a chamber which uses all rows of $\alpha$.  Condition (b) implies that this collection of rows is ``geometric'' and thus satisfies the hypotheses of the proposition.  Take $X$ to be a quotient which satisfies the conclusion.  By \cite[Theorem 14.2.13]{CLS2011} the space $X$ is a toric variety whose fan contains rays along the rows of $\alpha$, so all of these rays appear in its fan.  Since the entries of each row have no common factor the matrix $\alpha$ appears directly in the computation of $\Cl(X)$ \cite[Theorem 4.1.3]{CLS2011}, implying that $\deg$ is the degree function of torus-invariant divisors for some isomorphism $\Cl(X)\cong\Cl(S)$ \lauren{induced by a diagram chasing lemma or something about Gale duals?}.

    Finally \cite[Proposition 14.4.1]{CLS2011} implies that $X$ is semiprojective and its fan has support determined by the Gale dual of $\deg$.  Condition (c) implies that $\Eff(S)_\QQ$ is a pointed cone by Lemma~\ref{lem:positive-gradings}, so the support of the fan of $X$ is all of $N_\QQ$ by \cite[Proposition 14.3.2]{CLS2011}. \lauren{is strongly convex the same as pointed?} The space $X$ is semiprojective and proper, thus projective. \lauren{cite?}
  \end{proof}

\ifdefined\whole
\else
\bibliographystyle{plain}
\bibliography{Bibliography}
\end{document}
\fi
